\section*{Erd\H{o}s Problem 692: exactly one divisor in an interval}

\paragraph{Statement and status.}
For integers \(n<m\), let \(\delta_1(n,m)\) be the density of
integers having exactly one divisor in the open interval \((n,m)\).
The proposed unimodality in \(m\) is false. We reconstruct its
small counterexample and give a terminating procedure locating
the global maximum for every fixed \(n\). A useful explicit or
asymptotic formula for that location remains unresolved here.
The counterexample is attributed to Cambie,
\url{https://arxiv.org/abs/2501.10333};
the problem and its additional maximum-location question are at
\url{https://www.erdosproblems.com/692}.

\paragraph{Exact finite formula.}
Let \(D=\{n+1,\ldots,m-1\}\). The event is periodic modulo
\(\operatorname{lcm}(D)\), so the density exists. Inclusion--exclusion
with a marked divisor gives
\[
 \delta_1(n,m)=
 \sum_{\varnothing\ne S\subseteq D}
 \frac{(-1)^{|S|-1}|S|}{\operatorname{lcm}(S)}.                  \tag{1}
\]
To verify the coefficient, an integer divisible by exactly \(j\)
members of \(D\) contributes
\(\sum_{s=1}^j(-1)^{s-1}s\binom js\), which is \(1\) for \(j=1\)
and \(0\) for every other \(j\).

\paragraph{A strict decrease followed by a strict increase.}
For \(n=3\), the first relevant divisor sets are
\(\{4,5\}\), \(\{4,5,6\}\), and \(\{4,5,6,7\}\).
Formula (1) gives
\[
 \delta_1(3,6)=\frac14+\frac15-\frac2{20}=\frac7{20},
\]
\[
 \delta_1(3,7)=
 \left(\frac14+\frac15+\frac16\right)
 -2\left(\frac1{20}+\frac1{12}+\frac1{30}\right)
 +\frac3{60}=\frac13.
\]
The density of integers divisible by none of \(4,5,6\) is
\[
 1-\frac{37}{60}+\frac16-\frac1{60}=\frac8{15}.
\]
Divisibility by \(7\) is independent of the residue modulo \(60\).
Exactly one of \(4,5,6,7\) divides an integer precisely when it
has exactly one of the first three and is not divisible by \(7\),
or has none of the first three and is divisible by \(7\). Therefore
\[
 \delta_1(3,8)=\frac67\cdot\frac13+\frac17\cdot\frac8{15}
             =\frac{38}{105}.
\]
Since \(7/20>1/3<38/105\), unimodality fails.

\paragraph{The global maximum is attained and can be certified finitely.}
Fix \(n\geq1\). An integer with exactly one divisor in \((n,m)\)
has at most one prime divisor in that interval. Independence
over these primes gives the upper bound
\[
 \delta_1(n,m)\leq U_n(m):=
 \prod_{n<p<m}\left(1-\frac1p\right)
 \left(1+\sum_{n<p<m}\frac1{p-1}\right).                         \tag{2}
\]
Put \(S=\sum_{n<p<m}1/p\). Because
\(\sum_p(1/(p-1)-1/p)<\infty\), the right side is at most
\(e^{-S}(1+S+C)\), for an absolute constant \(C\).
Euler's divergence of the sum of reciprocal primes implies
\(S\longrightarrow\infty\), so \(U_n(m)\longrightarrow0\).
Moreover \(U_n(m)\) is nonincreasing: it is the probability
that at most one of an increasing collection of independent
prime-divisibility events occurs.

At \(m=n+2\), there is just the divisor \(n+1\), and the density
is \(1/(n+1)>0\). Enumerate primes until an integer cutoff \(M\)
satisfies \(U_n(M)<1/(n+1)\). This search terminates by (2).
For all \(m\geq M\), the density is smaller than the already
known value at \(n+2\). Evaluate (1) for the finitely many
integers \(n<m<M\), and select those with largest value.
This is a finite exact algorithm, with a certificate excluding
every later value of \(m\).

\paragraph{Remaining gap.}
The disproof of unimodality is complete. The procedure above
does not give an efficient bound, a simple formula, or an
asymptotic description of the maximizing location as \(n\) grows.
That further part of the question is unresolved in this attempt;
it should not be confused with the settled counterexample.

\paragraph{Completion estimate.}
\textbf{COMPLETION: 65\%}
